%! Author = sriadityavedantam
%! Date = 3/30/26

\section{Field Extensions}

Refer to Chapter 7.2 for all the theory on field extensions.
We will be continuing this on behalf of learning Galois Theory.

\begin{definition}
    Given $u\in\mathbb{E}\supseteq \mathbb{F}$, $\mathbb{F}(u)$ is the intersection of all subfields of $\mathbb{E}$ containing $\mathbb{F}$ and $u$.
    That is,
    \[
        \mathbb{F}(u)
        \coloneqq
        \bigcap_{\substack{\mathbb{E}_i\subseteq \mathbb{E}\\ \mathbb{F}\subseteq \mathbb{E}_i,\ u\in \mathbb{E}_i}} \mathbb{E}_i.
    \]
    It is the smallest field extension of $\mathbb{F}$ containing $u$.
    That is, if $\mathbb{K}$ is a field with
    \[
        \mathbb{F}\subseteq \mathbb{K}
        \quad\text{and}\quad
        u\in \mathbb{K},
    \]
    then
    \[
        \mathbb{F}(u)\subseteq \mathbb{K}.
    \]
\end{definition}

Thus if $u$ is algebraic over $\mathbb{F}$, then $\mathbb{F}[u]$ is a field extension of $\mathbb{F}$ with $u\in \mathbb{F}[u]$.
Hence
\[
    \mathbb{F}(u)\subseteq \mathbb{F}[u].
\]
Let $f(u)\in \mathbb{F}[u]$, where $f(x)\in \mathbb{F}[x]$.
Then it is always true that $f(u)\in \mathbb{F}(u)$, since $\mathbb{F}(u)$ must contain every polynomial in $u$.
Thus
\[
    \mathbb{F}[u]\subseteq \mathbb{F}(u).
\]
Therefore, when $u$ is algebraic over $\mathbb{F}$,
\[
    \mathbb{F}(u)=\mathbb{F}[u].
\]
Let $f(x)\in \mathbb{F}[x]$.
If $p(x)$ is an irreducible factor of $f(x)$, then
\[
    \frac{\mathbb{F}[x]}{(p(x))}
\]
is a field containing $\mathbb{F}$, really a copy consisting of classes of constant polynomials.
It also has a root of $p(x)$, namely the class of $x$.
For example, if
\[
    f(x)=x^2-2
\]
in $\mathbb{Q}[x]$, then
\[
    \mathbb{Q}\left[\sqrt{2}\right] \cong \frac{\mathbb{Q}[x]}{(x^2-2)}.
\]
Thus $\sqrt{2}$ corresponds to the class of $x$.
The map
\[
    \mathbb{Q}[x]\to \mathbb{Q}\left[\sqrt{2}\right]
\]
is defined by
\[
    f(x)\mapsto f\left(\sqrt{2}\right),
\]
so
\[
    x\mapsto \sqrt{2}.
\]
Then $p(\alpha)$ is equal to the class of $p(x)$.
But since we mod out by $(p(x))$, that is equal to
\[
    0\in \frac{\mathbb{F}[x]}{(p(x))}.
\]

\begin{theorem}{
    Let $p(x)\in \mathbb{F}_1[x]$ be irreducible.
    Let
    \[
        \mathbb{K}_1\coloneqq \frac{\mathbb{F}_1[x]}{(p(x))}.
    \]
    Then $\mathbb{K}_1$ is a field extension of $\mathbb{F}_1$ in which $p(x)$ has a root $\alpha\in \mathbb{K}_1$, where $\alpha$ is the class of $x$.
    Suppose
    \[
        \sigma:\mathbb{F}_1\to \mathbb{F}_2
    \]
    is a field isomorphism.
    If
    \[
        f(x)=a_0+a_{1}x+\cdots+a_{n}x^n\in \mathbb{F}_1[x],
    \]
    define
    \[
        \sigma f(x)\coloneqq \sigma(a_0)+\sigma(a_1)x+\cdots+\sigma(a_n)x^n\in \mathbb{F}_2[x].
    \]
    Then irreducibility is preserved.
    That is, if $p(x)$ is irreducible in $\mathbb{F}_1[x]$, then $\sigma p(x)$ is irreducible in $\mathbb{F}_2[x]$.
    Let $u\in \mathbb{E}_1$ be a root of $p(x)$.
    Suppose $v$ is a root of $\sigma p(x)\in \mathbb{F}_2[x]$.
    Then there exists an isomorphism
    \[
        \overline{\sigma}:\mathbb{F}_1(u)\to \mathbb{F}_2(v)
    \]
    extending $\sigma$, such that
    \[
        \overline{\sigma}|_{\mathbb{F}_1}=\sigma
        \quad\text{and}\quad
        \overline{\sigma}(u)=v.
    \]
}
\[
    \mathbb{F}_1(u)\cong \frac{\mathbb{F}_1[x]}{(p(x))}
    \cong \frac{\mathbb{F}_2[x]}{(\sigma p(x))}
    \cong \mathbb{F}_2(v).
\]
The first and third isomorphisms come from evaluation at $u$ and $v$, respectively.
The middle isomorphism is induced by $\sigma$ on coefficients.
Under these identifications, the class of $x$ in
\[
    \frac{\mathbb{F}_1[x]}{(p(x))}
\]
corresponds to the class of $x$ in
\[
    \frac{\mathbb{F}_2[x]}{(\sigma p(x))}
\]
and hence to $v$.
Thus we obtain an isomorphism
\[
    \overline{\sigma}:\mathbb{F}_1(u)\to \mathbb{F}_2(v)
\]
such that
\[
    \overline{\sigma}(u)=v.
\]
Also,
\[
    \overline{\sigma}(1_{\mathbb{F}_1})=1_{\mathbb{F}_2}.
\]
So $\overline{\sigma}$ is an extension of $\sigma$.
\end{theorem}

\begin{corollary}{
    The number of such extensions is at most the number of distinct roots of $\sigma p(x)$ in a chosen overfield containing $\mathbb{F}_2(v)$.
}
\end{corollary}

\begin{definition}[Splitting Field]
    A splitting field for $f(x)\in \mathbb{F}[x]$ is an extension field $\mathbb{E}$ of $\mathbb{F}$ such that in $\mathbb{E}[x]$ we have
    \[
        f(x)=c(x-a_1)\cdots(x-a_n),
        \qquad a_i\in \mathbb{E},
    \]
    and
    \[
        \mathbb{E}=\mathbb{F}(a_1,\ldots,a_n).
    \]
\end{definition}

\begin{theorem}{
    There exists a splitting field for every polynomial $f(x)\in \mathbb{F}[x]$.
}
We prove this by induction on $\deg(f)$.
If $\deg(f)=1$, then $f$ is already linear, so $\mathbb{F}$ itself is a splitting field.
Now assume $\deg(f)>1$.
Choose an irreducible factor $p(x)$ of $f(x)$.
Then
\[
    \mathbb{E}_1\coloneqq \frac{\mathbb{F}[x]}{(p(x))}
\]
is a field extension of $\mathbb{F}$ in which $p(x)$ has a root $u$.
Hence in $\mathbb{E}_1[x]$,
\[
    f(x)=(x-u)g(x)
\]
for some $g(x)$ with
\[
    \deg(g)<\deg(f).
\]
By the induction hypothesis, there exists a splitting field $\mathbb{E}$ for $g(x)$ over $\mathbb{E}_1$.
Then $\mathbb{E}$ is a splitting field for $f(x)$ over $\mathbb{F}$.
\end{theorem}

\begin{theorem}{
    Any two splitting fields for $f(x)$ over $\mathbb{F}$ are isomorphic over $\mathbb{F}$.
}
Suppose
\[
    \sigma:\mathbb{F}_1\isomto \mathbb{F}_2
\]
is a field isomorphism.
Let $f(x)\in \mathbb{F}_1[x]$.
Suppose $\mathbb{E}_1$ is a splitting field for $f(x)$ over $\mathbb{F}_1$, and $\mathbb{E}_2$ is a splitting field for $\sigma f(x)$ over $\mathbb{F}_2$.
We prove by induction on
\[
    [\mathbb{E}_1:\mathbb{F}_1]
\]
that $\sigma$ extends to an isomorphism
\[
    \overline{\sigma}:\mathbb{E}_1\to \mathbb{E}_2.
\]
If
\[
    [\mathbb{E}_1:\mathbb{F}_1]=1,
\]
then $\mathbb{E}_1=\mathbb{F}_1$, and the result is immediate.
Assume now
\[
    [\mathbb{E}_1:\mathbb{F}_1]>1.
\]
Let $p(x)$ be an irreducible factor of $f(x)$, and let $r\in \mathbb{E}_1$ be a root of $p(x)$.
If $v$ is any root of $\sigma p(x)$ in $\mathbb{E}_2$, then by the previous theorem there exists an extension
\[
    \sigma_1:\mathbb{F}_1(r)\isomto \mathbb{F}_2(v).
\]
Now $\mathbb{E}_1$ is a splitting field of $f(x)$ over $\mathbb{F}_1(r)$, and $\mathbb{E}_2$ is a splitting field of $\sigma_1 f(x)$ over $\mathbb{F}_2(v)$.
Since
\[
    [\mathbb{E}_1:\mathbb{F}_1(r)]<[\mathbb{E}_1:\mathbb{F}_1],
\]
the induction hypothesis applies.
Thus $\sigma_1$ extends to an isomorphism
\[
    \overline{\sigma}:\mathbb{E}_1\to \mathbb{E}_2.
\]
In particular, if
\[
    \mathbb{F}_1=\mathbb{F}_2=\mathbb{F}
\]
and $\sigma=\id_{\mathbb{F}}$, then any two splitting fields of $f(x)$ over $\mathbb{F}$ are isomorphic over $\mathbb{F}$.
Thus splitting fields are unique up to isomorphism.
\end{theorem}